\section{Baire空间}
\label{section:ch08-sec48}

就像本书中所引入的许多条件那样, 确定一个空间是 Baire 空间的条件也是"很不自然"的, 但是我们只能暂时先忍耐一下.

在本节中, 我们将定义 Baire 空间, 并指出两类重要的空间, 即完备度量空间和紧致的 Hausdorff 空间, 它们都包含在 Baire 空间类中. 然后再给出一些应用. 尽管这些应用并没有使 Baire 条件显得更自然一些, 但至少说明 Baire 空间是一个有用的工具. 事实上, 在分析学与拓扑学中它确实是一个十分有用而巧妙的工具.

\begin{definition}[空内部]
\label{definition:ch08-sec48-empty-interior}
我们知道, 空间 $X$ 的子集 $A$ 的内部, 是所有包含于 $A$ 的开集之并. 如果除了空集, $A$ 不包含 $X$ 的任何其他开集, 我们就说 $A$ 有\emph{空内部}(empty interior). 等价地, 如果 $A$ 的每一个点都是 $A$ 的补的一个极限点, 也就是说 $A$ 的补在 $X$ 中稠密, 我们就说 $A$ 有\emph{空内部}.
\end{definition}

\begin{example}
\label{example:ch08-sec48-1}
有理数集 $\mathbb{Q}$ 作为 $\mathbb{R}$ 的子集有空内部, 但闭区间 $[0,1]$ 内部非空. 区间 $[0,1]\times 0$ 作为平面 $\mathbb{R}^{2}$ 的子集有空内部, 子集 $\mathbb{Q}\times\mathbb{R}$ 也有空内部.
\end{example}

\begin{definition}[Baire空间]
\label{definition:ch08-sec48-baire}
空间 $X$ 称为一个\emph{Baire 空间}(Baire space), 如果它满足以下条件: 给定 $X$ 的闭集的任意可数族 $\{A_n\}$, 如果其中每一个集合在 $X$ 中有空内部, 那么它们的并 $\bigcup A_n$ 在 $X$ 中也有空内部.
\end{definition}

\begin{example}
\label{example:ch08-sec48-2}
有理数空间 $\mathbb{Q}$ 不是 Baire 空间, 因为 $\mathbb{Q}$ 的单点集为闭集, 并且在 $\mathbb{Q}$ 中有空内部, 但 $\mathbb{Q}$ 是它的单点子集的可数并.

另一方面, 空间 $\mathbb{Z}_{+}$ 是一个 Baire 空间. 因为 $\mathbb{Z}_{+}$ 的任意子集是开集, 所以除了空集外, $\mathbb{Z}_{+}$ 的任何子集都有非空内部, 这就是说, $\mathbb{Z}_{+}$ 绝对满足 Baire 条件.

更一般地, 作为完备度量空间, $\mathbb{R}$ 的任意闭子空间是一个 Baire 空间. 出人意料的是 $\mathbb{R}$ 中的无理数集也是一个 Baire 空间, 见习题 \cref{exercise:ch08-sec48-6}.
\end{example}

\begin{remark}
这个术语最初由 Baire(R. Baire)用在与"范畴"一词有关的概念之中. 空间 $X$ 的子集 $A$ 称为 $X$ 中的\emph{第一范畴集}, 如果它包含于具有空内部的 $X$ 的可数个闭集的并之中, 否则称 $A$ 是 $X$ 中的\emph{第二范畴集}. 使用这一术语, 我们可以说:
\begin{quote}
空间 $X$ 是一个 Baire 空间当且仅当 $X$ 中的任意非空开集是第二范畴集.
\end{quote}
在本书中, 我们不使用"第一范畴集"与"第二范畴集"这两个词.

上述定义是 Baire 空间的"闭集定义", 还有一个借助于开集的定义方式也经常使用. 参见以下引理.
\end{remark}

\begin{lemma}[开集形式]
\label{lemma:ch08-sec48-48-1}
$X$ 是一个 Baire 空间当且仅当 $X$ 中开集的任意可数族 $\{U_n\}$, 若每一个 $U_n$ 在 $X$ 中稠密, 则交 $\bigcap U_n$ 也在 $X$ 中稠密.
\end{lemma}

\begin{proof}
我们说过, 集合 $C$ 在 $X$ 中稠密是指 $\overline{C}=X$. 该引理可以立刻由下面两条推出:
\begin{enumerate}
\item $A$ 是 $X$ 中的一个闭集当且仅当 $X-A$ 是 $X$ 中的一个开集.
\item $B$ 在 $X$ 中有空内部当且仅当 $X-B$ 在 $X$ 中稠密.
\end{enumerate}
\end{proof}

有许多定理给出了确定一个空间是 Baire 空间的条件. 其中最重要的一个定理如下.

\begin{theorem}[Baire范畴定理]
\label{theorem:ch08-sec48-48-2}
若 $X$ 是一个紧致的 Hausdorff 空间或者是一个完备度量空间, 则 $X$ 是一个 Baire 空间.
\end{theorem}

\begin{proof}
给定 $X$ 的闭集的一个可数族 $\{A_n\}$, 其中每一个 $A_n$ 有空内部. 我们证明它们的并 $\bigcup A_n$ 在 $X$ 中也有空内部. 为此, 任意给定 $X$ 的一个非空开集 $U_0$, 我们必须找出一点 $x$, 它属于 $U_0$, 但不属于任何一个 $A_n$.

首先考虑集合 $A_1$. 根据假定, $A_1$ 不包含 $U_0$. 因此可以选出一点 $y$, 它属于 $U_0$, 但不在 $A_1$ 中. 由于 $A_1$ 是闭集, $X$ 是正则的, 所以可以选取 $y$ 的一个邻域 $U_1$, 使得
\[
\overline{U}_{1}\cap A_{1}=\varnothing,\qquad
\overline{U}_{1}\subset U_{0}.
\]
如果 $X$ 是一个度量空间, 还可以把 $U_1$ 选得足够小, 使其直径小于 $1$.

一般地, 给定非空开集 $U_{n-1}$, 可以选取 $U_{n-1}$ 的一个点, 使它不在闭集 $A_n$ 中, 然后选取这个点的一个邻域 $U_n$, 使得
\[
\overline{U}_{n}\cap A_{n}=\varnothing,\qquad
\overline{U}_{n}\subset U_{n-1},\qquad
\operatorname{diam}U_{n}<1/n,\quad\text{当 }X\text{ 是度量空间时}.
\]

我们断言 $\bigcap\overline{U}_{n}$ 非空. 根据这一结论立即可见定理成立. 因为若 $x$ 是 $\bigcap\overline{U}_{n}$ 的一个点, 则由于 $\overline{U}_{1}\subset U_{0}$, 可见 $x$ 属于 $U_{0}$. 由于对于每一个 $n$, $\overline{U}_{n}$ 与 $A_n$ 无交, 所以 $x$ 不属于 $A_n$.

证明 $\bigcap\overline{U}_{n}$ 非空可分成两部分, 取决于 $X$ 是紧致的 Hausdorff 空间还是完备度量空间. 如果 $X$ 是一个紧致的 Hausdorff 空间, 考虑 $X$ 的非空子集的一个套序列 $\overline{U}_{1}\supset\overline{U}_{2}\supset\cdots$. 族 $\{\overline{U}_{n}\}$ 满足有限交性质, 由于 $X$ 是紧致的, $\bigcap\overline{U}_{n}$ 必然非空.

如果 $X$ 是一个完备度量空间, 则可以应用以下引理完成定理的证明.
\end{proof}

\begin{lemma}
\label{lemma:ch08-sec48-48-3}
设 $C_{1}\supset C_{2}\supset\cdots$ 是完备度量空间 $X$ 中的一个非空闭集的套序列. 若 $\operatorname{diam}C_{n}\to 0$, 则 $\bigcap C_{n}\neq\varnothing$.
\end{lemma}

\begin{proof}
这个引理在第 43 节的习题中曾出现过, 这里给出证明. 对于每一个 $n$, 选取 $x_{n}\in C_{n}$. 因为对于 $n,m\geqslant N$, 有 $x_{n},x_{m}\in C_{N}$, 并且对任意给定的 $\varepsilon>0$, 可以选取充分大的 $N$, 使得 $C_{N}$ 的直径小于 $\varepsilon$, 所以序列 $(x_n)$ 是一个 Cauchy 序列. 假定 $(x_n)$ 收敛于 $x$, 则对于任意给定的 $k$, 子序列 $x_k,x_{k+1},\cdots$ 也收敛于 $x$. 于是, $x$ 必定属于 $\overline{C}_{k}=C_{k}$, 因此 $x\in\bigcap C_{k}$.
\end{proof}

这里有 Baire 空间理论的一个应用, 本节后面我们还将给出更多的应用. 这个应用有些意思但不深奥, 它涉及的是一个学生可能会问到的关于收敛的连续函数序列的问题.

设 $f_n\colon[0,1]\to\mathbb{R}$ 是连续函数的一个序列, 对于每一个 $x\in[0,1]$, 有 $f_n(x)\to f(x)$. 已有例子说明极限函数 $f$ 未必是连续的. 但人们想知道的是 $f$ 到底是如何不连续的呢? 例如, 它是处处不连续的吗? 答案是否定的. 我们将要证明 $f$ 必定在 $[0,1]$ 的无限多个点处连续. 事实上, 所有使得 $f$ 连续的点构成的集合在 $[0,1]$ 中稠密.

为证明这个结论, 需要以下引理.

\begin{lemma}
\label{lemma:ch08-sec48-48-4}
Baire 空间 $X$ 的任何一个开子空间 $Y$ 是一个 Baire 空间.
\end{lemma}

\begin{proof}
设 $A_n$ 是 $Y$ 中有空内部的闭集的一个可数族, 下面证明 $\bigcup A_n$ 在 $Y$ 中有空内部.

设 $\overline{A}_{n}$ 是 $A_n$ 在 $X$ 中的闭包, 则 $\overline{A}_{n}\cap Y=A_n$. 于是, $\overline{A}_{n}$ 在 $X$ 中有空内部. 因为若 $U$ 是 $X$ 中包含于 $\overline{A}_{n}$ 的一个非空开集, 则 $U$ 必定与 $A_n$ 相交, 这导致 $U\cap Y$ 是 $Y$ 中包含在 $A_n$ 中的一个非空开集, 与假设矛盾.

若集合 $A_n$ 的并包含着 $Y$ 中的一个非空开集 $W$, 则 $\overline{A}_{n}$ 的并也包含着 $W$. 由于 $Y$ 在 $X$ 中是开的, 所以 $W$ 在 $X$ 中也是开的. 但是每一个集合 $\overline{A}_{n}$ 在 $X$ 中有空内部, 这与 $X$ 是一个 Baire 空间矛盾.
\end{proof}

\begin{theorem}
\label{theorem:ch08-sec48-48-5}
设 $X$ 是一个空间, $(Y,d)$ 是一个度量空间. 设 $f_n\colon X\to Y$ 是一个连续函数序列, 使得对于所有的 $x\in X$, 有 $f_n(x)\to f(x)$, 其中 $f\colon X\to Y$. 如果 $X$ 是一个 Baire 空间, 那么使得 $f$ 连续的点构成的集合在 $X$ 中稠密.
\end{theorem}

\begin{proof}
给定正整数 $N$ 和 $\varepsilon>0$, 定义
\[
A_{N}(\varepsilon)=\{\,x\mid d(f_{n}(x),f_{m}(x))\leqslant\varepsilon,\ \text{对于所有 }n,m\geqslant N\,\}.
\]
注意, $A_N(\varepsilon)$ 在 $X$ 中是闭的. 因为根据 $f_n$ 和 $f_m$ 的连续性可见, 满足 $d(f_n(x),f_m(x))\leqslant\varepsilon$ 的点的集合是 $X$ 的一个闭集, 并且对于所有 $n,m\geqslant N$, $A_N(\varepsilon)$ 是这些集合的交, 即
\[
A_{N}(\varepsilon)=\bigcap_{n,m\geqslant N}\{\,x\mid d(f_{n}(x),f_{m}(x))\leqslant\varepsilon\,\}.
\]

对于固定的 $\varepsilon$, 考虑集合 $A_1(\varepsilon)\subset A_2(\varepsilon)\subset\cdots$. 这些集合的并为 $X$. 因为任意给定 $x_0\in X$, 由于 $f_n(x_0)\to f(x_0)$ 可得序列 $f_n(x_0)$ 是一个 Cauchy 序列, 因此对于某一个 $N$, 有 $x_0\in A_N(\varepsilon)$.

现在, 设
\[
U(\varepsilon)=\bigcup_{N\in\mathbb{Z}_{+}}\operatorname{Int}A_{N}(\varepsilon).
\]
我们来证明两件事情:
\begin{enumerate}
\item $U(\varepsilon)$ 是 $X$ 中的一个稠密开集.
\item 函数 $f$ 在集合
\[
C=U(1)\cap U(1/2)\cap U(1/3)\cap\cdots
\]
的每一点处都连续.
\end{enumerate}
由此推出, 本定理在 $X$ 是 Baire 空间的假设下成立.

为了证明 $U(\varepsilon)$ 在 $X$ 中是稠密的, 需要证明对于 $X$ 的任意一个非空开集 $V$, 存在 $N$ 使得 $V\cap\operatorname{Int}A_N(\varepsilon)$ 是非空的. 为此, 我们先注意对每一个 $N$, $V\cap A_N(\varepsilon)$ 在 $V$ 中是闭的. 因为根据前述引理可见 $V$ 是一个 Baire 空间, 所以至少存在一个这种集合, 不妨设为 $V\cap A_M(\varepsilon)$, 必包含 $V$ 中的一个非空开集 $W$. 因为 $V$ 在 $X$ 中是开的, 所以 $W$ 也是 $X$ 中的开集. 因此, 它包含于 $\operatorname{Int}A_M(\varepsilon)$.

现在我们来证明, 若 $x_0\in C$, 则 $f$ 在点 $x_0$ 处连续. 给定 $\varepsilon>0$, 我们来找 $x_0$ 的一个邻域 $W$, 使得对于 $x\in W$, $d(f(x),f(x_0))<\varepsilon$.

首先, 选取 $k$ 使得 $1/k<\varepsilon/3$. 因为 $x_0\in C$, 所以 $x_0\in U(1/k)$. 因此存在某一个 $N$, 使得 $x_0\in\operatorname{Int}A_N(1/k)$. 最后, 由 $f_N$ 的连续性可以找到 $x_0$ 的一个邻域 $W$ 包含于 $A_N(1/k)$, 使得
\[
d(f_{N}(x),f_{N}(x_{0}))<\varepsilon/3,\quad\text{对于 }x\in W. \tag{$*$}
\]
根据 $W\subset A_N(1/k)$ 可见,
\[
d(f_{n}(x),f_{N}(x))\leqslant 1/k,\quad\text{对于 }n\geqslant N\text{ 和 }x\in W.
\]
令 $n\to\infty$, 我们得到不等式
\[
d(f(x),f_{N}(x))\leqslant 1/k<\varepsilon/3,\quad\text{对于 }x\in W. \tag{$**$}
\]
特别地, 因为 $x_0\in W$, 我们有
\[
d(f(x_{0}),f_{N}(x_{0}))<\varepsilon/3. \tag{$***$}
\]
对于 $(*)$、$(**)$ 和 $(***)$ 应用三角不等式, 便得到结论.
\end{proof}

\subsection*{习题}

\begin{exercise}
设 $X$ 等于可数并 $\bigcup B_n$. 证明: 如果 $X$ 是一个非空的 Baire 空间, 则至少存在一个集合 $\overline{B}_n$ 有非空内部.
\end{exercise}

\begin{exercise}
根据 Baire 范畴定理可见 $\mathbb{R}$ 不能写成可数个具有空内部的闭子集的并. 证明如果不要求这些集合是闭的, 那么上面的结论不成立.
\end{exercise}

\begin{exercise}
证明每一个局部紧致的 Hausdorff 空间是 Baire 空间.
\end{exercise}

\begin{exercise}
证明: 若 $X$ 中的每一个点 $x$ 都有一个邻域是 Baire 空间, 则 $X$ 是 Baire 空间. [提示: 应用 Baire 条件的开集形式.]
\end{exercise}

\begin{exercise}
证明: 若 $Y$ 为 $X$ 中的一个 $G_\delta$ 集, 并且 $X$ 是紧致的 Hausdorff 空间或完备度量空间, 则 $Y$ 在子空间拓扑下是一个 Baire 空间. [提示: 假定 $Y=\bigcap W_n$, 其中 $W_n$ 在 $X$ 中是开的, 并假定 $B_n$ 是 $Y$ 中具有空内部的一个闭集. 给定 $X$ 的一个开集 $U_0$, 使得 $U_0\cap Y\neq\varnothing$, 找出 $X$ 中一个开集序列 $U_n$, 满足 $U_n\cap Y\neq\varnothing$, 并使得
\[
\overline{U}_{n}\subset U_{n-1},\qquad
\overline{U}_{n}\cap\overline{B}_{n}=\varnothing,\qquad
\operatorname{diam}U_{n}<1/n,\quad\text{当 }X\text{ 是度量空间时},\qquad
\overline{U}_{n}\subset W_{n}.
\]
]
\end{exercise}

\begin{exercise}
\\label{exercise:ch08-sec48-6}
证明无理数集是一个 Baire 空间.
\end{exercise}

\begin{exercise}
\textbf{定理} 若 $D$ 是 $\mathbb{R}$ 的一个可数稠密子集, 则没有函数 $f\colon\mathbb{R}\to\mathbb{R}$ 恰在 $D$ 的所有点处连续.
\begin{enumerate}
\item 证明: 如果 $f\colon\mathbb{R}\to\mathbb{R}$, 那么使得 $f$ 连续的所有点构成的集合 $C$ 是 $\mathbb{R}$ 中的一个 $G_\delta$ 集. [提示: 设 $U_n$ 是 $\mathbb{R}$ 中所有满足 $\operatorname{diam}f(U)<1/n$ 的开集 $U$ 的并, 证明 $C=\bigcap U_n$.]
\item 证明 $D$ 不是 $\mathbb{R}$ 中的一个 $G_\delta$ 集. [提示: 设 $D=\bigcap W_n$, 其中 $W_n$ 在 $\mathbb{R}$ 中是开的. 对于 $d\in D$, 令 $V_d=\mathbb{R}-\{d\}$. 证明 $W_n$ 和 $V_d$ 都在 $\mathbb{R}$ 中稠密.]
\end{enumerate}
\end{exercise}

\begin{exercise}
设 $f_n$ 是一个连续函数序列, $f_n\colon\mathbb{R}\to\mathbb{R}$, 使得对任意 $x\in\mathbb{R}$, 有 $f_n(x)\to f(x)$. 证明 $f$ 在 $\mathbb{R}$ 中的不可数多个点处连续.
\end{exercise}

\begin{exercise}
设 $g\colon\mathbb{Z}_{+}\to\mathbb{Q}$ 是一个一一映射, 令 $x_n=g(n)$. 定义 $f\colon\mathbb{R}\to\mathbb{R}$ 如下:
\[
f(x_{n})=1/n,\quad\text{对于 }x_{n}\in\mathbb{Q},\qquad
f(x)=0,\quad\text{对于 }x\notin\mathbb{Q}.
\]
证明 $f$ 在每一个无理数处连续, 在每一个有理数处不连续. 你能找到一个连续函数序列 $f_n$ 收敛到 $f$ 吗?
\end{exercise}

\begin{exercise}
证明以下结论:

\textbf{定理(一致有界原理)} 设 $X$ 是一个完备度量空间, $\mathcal{F}$ 是 $\mathcal{C}(X,\mathbb{R})$ 的一个子集, 使得对于每一个 $a\in X$, 集合
\[
\mathcal{F}_{a}=\{\,f(a)\mid f\in\mathcal{F}\,\}
\]
有界, 则存在 $X$ 中的一个非空开集 $U$, 使得 $\mathcal{F}$ 中的函数在 $U$ 上一致有界, 即存在一个数 $M$, 使得对于所有的 $x\in U$ 及所有 $f\in\mathcal{F}$, 有 $|f(x)|\leqslant M$. [提示: 令 $A_N=\{\,x;\ |f(x)|\leqslant N,\ \text{对于所有 }f\in\mathcal{F}\,\}$.]
\end{exercise}

\begin{exercise}
判定 $\mathbb{R}_\ell$ 是否为 Baire 空间.
\end{exercise}

\begin{exercise}
证明 $\mathbb{R}^\omega$ 在箱拓扑、积拓扑以及一致拓扑下都是 Baire 空间.
\end{exercise}

\begin{exercise}
设 $X$ 是一个拓扑空间, $Y$ 是一个完备度量空间. 证明 $\mathcal{C}(X,Y)$ 在细拓扑(参见第 46 节的习题 11)下是一个 Baire 空间. [提示: 给定基元素 $B(f_i,\delta_i)$ 使得 $\delta_1\leqslant 1$, $\delta_{i+1}\leqslant\delta_i/3$, 并且 $f_{i+1}\in B(f_i,\delta_i/3)$, 证明 $\bigcap B(f_i,\delta_i)\neq\varnothing$.]
\end{exercise}
